\begin{exo}{}
	Les vétérinaires donnent parfois le tableau de correspondance entre l'âge des chats et l'équivalent en âge humain ci-dessous.


	\begin{center}
		\renewcommand{\arraystretch}{1.5}
		\begin{tabular}{|c|c|c|c|c|c|c|} \hline
			\cellcolor{gray!30}\footnotesize\textbf{Âge du chat (en mois)} & $0,5$ & $1$  & $2$  & $6$  & $12$ & $16$ \\ \hline
			\cellcolor{gray!30}\footnotesize\textbf{Âge humain (en année)} & $10$  & $18$ & $26$ & $42$ & $70$ & $94$ \\ \hline
		\end{tabular}
	\end{center}


	On note $c$ l'âge du chat en mois et $H(c)$ l'âge humain équivalent en année.

	\begin{enumerate}
		\item Dans un repère orthogonal, tracer une courbe représentant la fonction $H$ sur $\intff{0}{16}$.
		\item Les deux âges sont-ils proportionnels ? Justifier.\\
		      \textit{\emph{Aide :} Quelle est la représentation graphique qui modélise une situation de proportionnalité ?}
		\item Préciser l'image de $3$ et interpréter la réponse.
		\item Donner un antécédent de $60$ et interpréter la réponse.
	\end{enumerate}
\end{exo}
